Bagian ini memuat uraian singkat (tidak lebih dari 250 kata) tetapi padat dan jelas serta memberikan gambaran menyeluruh tentang isi tesis. Intisari memuat apa dan mengapa penelitian dikerjakan, bagaimana dikerjakan, dan apa hasil penting yang diperoleh dari penelitian. Intisari ditulis dalam satu atau beberapa paragraf.

Contoh:

Deteksi objek secara \textit{real-time} merupakan komponen krusial dalam berbagai aplikasi \textit{computer vision} seperti kendaraan otonom dan sistem pengawasan cerdas. Arsitektur YOLOv8 telah menunjukkan performa deteksi yang unggul, namun ukuran model yang besar menjadi kendala utama untuk \textit{deployment} pada perangkat \textit{edge} dengan sumber daya komputasi terbatas. Penelitian ini bertujuan untuk mengembangkan model YOLOv8 yang dioptimasi melalui kombinasi \textit{structured pruning} dan \textit{post-training quantization} agar dapat dijalankan secara efisien pada perangkat NVIDIA Jetson Nano. Metode yang digunakan meliputi pelatihan model \textit{baseline} pada dataset COCO 2017, penerapan \textit{structured pruning} dengan rasio 30\%, 50\%, dan 70\% pada \textit{backbone} dan \textit{neck} model, dilanjutkan dengan \textit{fine-tuning} dan konversi ke format INT8 menggunakan TensorRT. Hasil eksperimen menunjukkan bahwa kombinasi \textit{pruning} 50\% dan \textit{quantization} INT8 mampu mengurangi ukuran model sebesar 73\% dan meningkatkan kecepatan inferensi hingga 2,8$\times$ (dari 12 FPS menjadi 34 FPS pada Jetson Nano) dengan penurunan mAP@0.5 hanya sebesar 2,1\% dibandingkan model \textit{baseline}. Hasil ini menunjukkan bahwa teknik optimasi yang diusulkan efektif untuk menghasilkan model deteksi objek yang ringan dan cepat tanpa mengorbankan akurasi secara signifikan.

\vspace{2ex}
\noindent \textbf{Kata Kunci: deteksi objek, YOLOv8, \textit{pruning}, \textit{quantization}, perangkat \textit{edge}}
